\chapter{1992年全国卷（理）}
\section{单选题}

\begin{problem}
\(\frac{\log_8 9}{\log_2 3}\) 的值是（\quad）

\choices
    {\(\frac{2}{3}\)}
    {\(1\)}
    {\(\frac{3}{2}\)}
    {\(2\)}

\begin{answer}
A
\end{answer}

\begin{solution}
由换底公式，且 \(\log_2 3\neq0\)，有
\(\log_8 9=\frac{\log_2 9}{\log_2 8}=\frac{2\log_2 3}{3}\). 因而
\(\frac{\log_8 9}{\log_2 3}=\frac{2}{3}\)，对应选项 A.
\end{solution}
\end{problem}

\begin{problem}
如果函数 \(y=\sin(\omega x)\cos(\omega x)\) 的最小正周期是 \(4\pi\)，那么常数 \(\omega\) 为（\quad）

\choices
    {\(4\)}
    {\(2\)}
    {\(\frac{1}{2}\)}
    {\(\frac{1}{4}\)}

\begin{answer}
D
\end{answer}

\begin{solution}
由二倍角公式，\(\sin(\omega x)\cos(\omega x)=\frac{1}{2}\sin(2\omega x)\). 当 \(\omega\neq0\) 时，\(\sin(2\omega x)\) 的最小正周期为 \(\frac{2\pi}{|2\omega|}=\frac{\pi}{|\omega|}\). 因此
\(\frac{\pi}{|\omega|}=4\pi\)，得 \(|\omega|=\frac{1}{4}\). 按选项中的通常约定取正的参数值，\(\omega=\frac{1}{4}\)，故选 D.
\end{solution}
\end{problem}

\begin{problem}
极坐标方程分别是 \(\rho=\cos\theta\) 和 \(\rho=\sin\theta\) 的两个圆的圆心距是（\quad）

\choices
    {\(2\)}
    {\(\sqrt{2}\)}
    {\(1\)}
    {\(\frac{\sqrt{2}}{2}\)}

\begin{answer}
D
\end{answer}

\begin{solution}
由 \(x=\rho\cos\theta\)、\(y=\rho\sin\theta\) 及 \(\rho^2=x^2+y^2\)，方程 \(\rho=\cos\theta\) 可化为
\(x^2+y^2=\rho^2=\rho\cos\theta=x\)，即
\(\left(x-\frac{1}{2}\right)^2+y^2=\frac{1}{4}\). 它的圆心为 \(\left(\frac{1}{2},0\right)\).

同理，方程 \(\rho=\sin\theta\) 化为 \(x^2+y^2=y\)，即
\(x^2+\left(y-\frac{1}{2}\right)^2=\frac{1}{4}\)，圆心为 \(\left(0,\frac{1}{2}\right)\). 所以圆心距为
\(\sqrt{\left(\frac{1}{2}\right)^2+\left(-\frac{1}{2}\right)^2}=\frac{\sqrt{2}}{2}\)，故选 D.
\end{solution}
\end{problem}

\begin{problem}
\(\sin 4x\cos 5x=-\cos 4x\sin 5x\) 的一个解是（\quad）

\choices
    {\(10^{\circ}\)}
    {\(20^{\circ}\)}
    {\(50^{\circ}\)}
    {\(70^{\circ}\)}

\begin{answer}
B
\end{answer}

\begin{solution}
将右边移到左边，并用正弦加法公式，得
\(\sin4x\cos5x+\cos4x\sin5x=\sin(4x+5x)=\sin9x=0\). 因而 \(9x=180^{\circ}k\)，即 \(x=20^{\circ}k\)，其中 \(k\in\Z\). 四个选项中只有 \(20^{\circ}\) 对应整数 \(k=1\)，故选 B.
\end{solution}
\end{problem}

\begin{problem}
已知轴截面是正方形的圆柱的高与球的直径相等，则圆柱的全面积与球的表面积的比是（\quad）

\choices
    {\(6:5\)}
    {\(5:4\)}
    {\(4:3\)}
    {\(3:2\)}

\begin{answer}
D
\end{answer}

\begin{solution}
设圆柱底面半径为 \(r\)，高为 \(h\). 圆柱的轴截面是长为 \(h\)、宽为 \(2r\) 的矩形，所以正方形条件给出 \(h=2r\). 又因为圆柱的高等于球的直径，球半径为 \(R=\frac h2=r\).

圆柱的全面积为 \(2\pi r^2+2\pi rh=2\pi r^2+4\pi r^2=6\pi r^2\)，球的表面积为 \(4\pi R^2=4\pi r^2\). 因而所求比为 \(6\pi r^2:4\pi r^2=3:2\)，故选 D.
\end{solution}
\end{problem}

\begin{problem}
如图，图中曲线是幂函数 \(y=x^n\) 在第一象限的图象. 已知 \(n\) 取 \(\pm 2\)，\(\pm \frac{1}{2}\) 四个值，则相应于曲线 \(C_1\)，\(C_2\)，\(C_3\)，\(C_4\) 的 \(n\) 依次为（\quad）

\bitmapfigure[max width=0.28\linewidth]{img/1992/national_science/q6_fig1.png}

\choices
    {\(-2, -\frac{1}{2}, \frac{1}{2}, 2\)}
    {\(2\)，\(\frac{1}{2}\)，\(-\frac{1}{2}\)，\(-2\)}
    {\(-\frac{1}{2}\)，\(-2\)，\(2\)，\(\frac{1}{2}\)}
    {\(2\)，\(\frac{1}{2}\)，\(-2\)，\(-\frac{1}{2}\)}

\begin{answer}
B
\end{answer}

\begin{solution}
四条曲线都经过 \((1,1)\). 取图中 \(x>1\) 的部分，则
\(x^2>x^{1/2}>x^{-1/2}>x^{-2}\). 因此从上到下，曲线依次对应指数 \(2\)、\(\frac12\)、\(-\frac12\)、\(-2\). 这与图中的 \(C_1\)，\(C_2\)，\(C_3\)，\(C_4\) 依次相符，故选 B.
\end{solution}
\end{problem}

\begin{problem}
若 \(\log_a 2<\log_b 2<0\)，则（\quad）

\choices
    {\(0<a<b<1\)}
    {\(0<b<a<1\)}
    {\(a>b>1\)}
    {\(b>a>1\)}

\begin{answer}
B
\end{answer}

\begin{solution}
对数有意义要求底数为正且不等于 \(1\). 因为 \(\log_a2<0\)、\(\log_b2<0\)，而 \(\ln2>0\)，所以 \(\ln a<0\)、\(\ln b<0\)，从而 \(0<a<1\)、\(0<b<1\).

令 \(A=-\ln a>0\)、\(B=-\ln b>0\)，则
\(\log_a2=-\frac{\ln2}{A}\)、\(\log_b2=-\frac{\ln2}{B}\). 由题设不等式，除以负数 \(-\ln2\) 时不等号反向，得到 \(\frac1A>\frac1B\)，所以 \(A<B\). 这等价于 \(\ln a>\ln b\)，即 \(a>b\). 因而 \(0<b<a<1\)，故选 B.
\end{solution}
\end{problem}

\begin{problem}
直线 \(\begin{cases}
x=t\cdot \sin 20^{\circ}+3,\\
y=-t\cdot \cos 20^{\circ}
\end{cases}\)（\(t\) 为参数）的倾斜角是（\quad）

\choices
    {\(20^{\circ}\)}
    {\(70^{\circ}\)}
    {\(110^{\circ}\)}
    {\(160^{\circ}\)}

\begin{answer}
C
\end{answer}

\begin{solution}
参数方程给出的方向向量为 \(\left(\sin20^{\circ},-\cos20^{\circ}\right)\)，所以直线斜率为
\(\frac{-\cos20^{\circ}}{\sin20^{\circ}}=-\cot20^{\circ}=\tan110^{\circ}\). 由于倾斜角取 \([0^{\circ},180^{\circ})\) 内的角，故倾斜角为 \(110^{\circ}\)，选 C.
\end{solution}
\end{problem}

\begin{problem}
在四棱锥的四个侧面中，直角三角形最多可有（\quad）

\choices
    {\(1\) 个}
    {\(2\) 个}
    {\(3\) 个}
    {\(4\) 个}

\begin{answer}
D
\end{answer}

\begin{solution}
四棱锥只有四个侧面，所以直角三角形的个数不可能超过 \(4\). 下面验证上界可以达到：取
\(A=(0,0,0)\)、\(B=(1,0,0)\)、\(C=(1,1,0)\)、\(D=(0,1,0)\)，并取顶点 \(S=(0,0,h)\)，其中 \(h>0\). 则
\(\overrightarrow{AS}\cdot\overrightarrow{AB}=0\)、\(\overrightarrow{AS}\cdot\overrightarrow{AD}=0\)，所以侧面 \(SAB\)、\(SAD\) 都是直角三角形；又
\(\overrightarrow{BS}=(-1,0,h)\)、\(\overrightarrow{BC}=(0,1,0)\)，以及 \(\overrightarrow{DS}=(0,-1,h)\)、\(\overrightarrow{DC}=(1,0,0)\)，两组内积都为零，所以侧面 \(SBC\)、\(SCD\) 也都是直角三角形. 因而最大个数为 \(4\)，故选 D.
\end{solution}
\end{problem}

\begin{problem}
圆心在抛物线 \(y^2=2x\) 上，且与 \(x\) 轴和该抛物线的准线都相切的一个圆的方程是（\quad）

\choices
    {\(x^2+y^2-x-2y-\frac{1}{4}=0\)}
    {\(x^2+y^2+x-2y+1=0\)}
    {\(x^2+y^2-x-2y+1=0\)}
    {\(x^2+y^2-x-2y+\frac{1}{4}=0\)}

\begin{answer}
D
\end{answer}

\begin{solution}
设圆心为 \((u,v)\)，半径为 \(r\). 由圆心在 \(y^2=2x\) 上，得 \(v^2=2u\)，即 \(u=\frac{v^2}{2}\geqslant0\). 该抛物线的准线为 \(x=-\frac12\). 圆与 \(x\) 轴相切，故 \(r=|v|\)；圆与准线相切，故
\(r=\left|u+\frac12\right|=u+\frac12\). 于是
\( |v|=\frac{v^2}{2}+\frac12\). 令 \(s=|v|\)，则 \(s=\frac{s^2}{2}+\frac12\)，即 \((s-1)^2=0\)，所以 \(s=1\)、\(u=\frac12\)、\(r=1\).

因此圆心可能为 \(\left(\frac12,1\right)\) 或 \(\left(\frac12,-1\right)\). 前者的方程为
\(\left(x-\frac12\right)^2+(y-1)^2=1\)，化为 \(x^2+y^2-x-2y+\frac14=0\)，正是选项 D；后者关于 \(x\) 轴对称且不在给出的选项中. 故选 D.
\end{solution}
\end{problem}

\begin{problem}
在 \((x^2+3x+2)^5\) 的展开式中 \(x\) 的系数为（\quad）

\choices
    {\(160\)}
    {\(240\)}
    {\(360\)}
    {\(800\)}

\begin{answer}
B
\end{answer}

\begin{solution}
把 \((x^2+3x+2)^5\) 看成五个相同因式的乘积. 要得到一次项 \(x\)，每个因式至多取一次一次项，且不能取 \(x^2\)；因此必须从五个因式中选一个取 \(3x\)，其余四个因式都取常数项 \(2\). 所以 \(x\) 的系数为
\(\mathrm C_5^1\cdot3\cdot2^4=240\)，故选 B.
\end{solution}
\end{problem}

\begin{problem}
若 \(0<a<1\)，在 \([0,2\pi]\) 上满足 \(\sin x\geqslant a\) 的 \(x\) 的范围是（\quad）

\choices
    {\([0,\arcsin a]\)}
    {\([\arcsin a,\pi-\arcsin a]\)}
    {\([\pi-\arcsin a,\pi]\)}
    {\(\left[\arcsin a,\frac{\pi}{2}+\arcsin a\right]\)}

\begin{answer}
B
\end{answer}

\begin{solution}
令 \(\alpha=\arcsin a\). 因为 \(0<a<1\)，所以 \(0<\alpha<\frac\pi2\)，且在 \([0,\pi]\) 上有 \(\sin\alpha=a\)、\(\sin(\pi-\alpha)=a\). 正弦函数在 \([0,\pi]\) 上先增后减，故 \(\sin x\geqslant a\) 等价于 \(\alpha\leqslant x\leqslant\pi-\alpha\). 在 \([\pi,2\pi]\) 上 \(\sin x\leqslant0<a\)，没有新的解. 因而 \(x\in[\arcsin a,\pi-\arcsin a]\)，故选 B.
\end{solution}
\end{problem}

\begin{problem}
已知直线 \(l_1\) 和 \(l_2\) 夹角的平分线为 \(y=x\)，如果 \(l_1\) 的方程是 \(ax+by+c=0\)（\(ab>0\)），那么 \(l_2\) 的方程是（\quad）

\choices
    {\(bx+ay+c=0\)}
    {\(ax-by+c=0\)}
    {\(bx+ay-c=0\)}
    {\(bx-ay+c=0\)}

\begin{answer}
A
\end{answer}

\begin{solution}
关于直线 \(y=x\) 对称会交换坐标 \(x\)，\(y\). 因而直线 \(l_1:ax+by+c=0\) 的对称直线为
\(l_2:ay+bx+c=0\)，即 \(bx+ay+c=0\). 这两条直线关于 \(y=x\) 对称，所以 \(y=x\) 是它们的一条夹角平分线.

也可直接验证：两直线的法向量长度同为 \(\sqrt{a^2+b^2}\)，而点到两直线的距离相等时满足
\(\left|ax+by+c\right|=\left|bx+ay+c\right|\). 其中一个等距轨迹分支给出
\((a-b)(x-y)=0\). 两直线不重合时 \(a\neq b\)，故该分支即为 \(y=x\). 因此选项 A 正确，故选 A.
\end{solution}
\end{problem}

\begin{problem}
如图，在棱长为 \(1\) 的正方体 \(ABCD-A_1B_1C_1D_1\) 中，\(M\) 和 \(N\) 分别为 \(A_1B_1\) 和 \(BB_1\) 的中点，那么直线 \(AM\) 与 \(CN\) 所成角的余弦值是（\quad）

\bitmapfigure[max width=5.3cm]{img/1992/national_science/q14_fig1.png}

\choices
    {\(\frac{\sqrt{3}}{2}\)}
    {\(\frac{\sqrt{10}}{10}\)}
    {\(\frac{3}{5}\)}
    {\(\frac{2}{5}\)}

\begin{answer}
D
\end{answer}

\begin{solution}
建立直角坐标系，取 \(A=(0,0,0)\)、\(B=(1,0,0)\)、\(C=(1,1,0)\)、\(A_1=(0,0,1)\). 因为 \(M\)，\(N\) 分别为中点，所以
\(M=\left(\frac12,0,1\right)\)、\(N=\left(1,0,\frac12\right)\). 于是
\(\overrightarrow{AM}=\left(\frac12,0,1\right)\)、\(\overrightarrow{CN}=\left(0,-1,\frac12\right)\).

两向量的内积为 \(\overrightarrow{AM}\cdot\overrightarrow{CN}=\frac12\)，且
\(\lvert\overrightarrow{AM}\rvert=\lvert\overrightarrow{CN}\rvert=\sqrt{\frac14+1}=\frac{\sqrt5}{2}\). 因此两直线所成锐角的余弦为
\(\frac{\left|\overrightarrow{AM}\cdot\overrightarrow{CN}\right|}{\lvert\overrightarrow{AM}\rvert\lvert\overrightarrow{CN}\rvert}=\frac{\frac12}{\frac54}=\frac25\)，故选 D.
\end{solution}
\end{problem}

\begin{problem}
已知复数 \(z\) 的模为 \(2\)，则 \(|z-\i|\) 的最大值为（\quad）

\choices
    {\(1\)}
    {\(2\)}
    {\(\sqrt{5}\)}
    {\(3\)}

\begin{answer}
D
\end{answer}

\begin{solution}
由三角不等式，\(|z-\i|=|z+(-\i)|\leqslant|z|+|\i|=2+1=3\). 还需验证等号可以取得：取 \(z=-2\i\)，则 \(|z|=2\)，且 \(z\) 与 \(-\i\) 同向，所以
\(|z-\i|=|-3\i|=3\). 因而最大值为 \(3\)，故选 D.
\end{solution}
\end{problem}

\begin{problem}
函数 \(y=\frac{\e^x-\e^{-x}}{2}\) 的反函数（\quad）

\choices
    {是奇函数，它在 \((0,+\infty)\) 上是减函数}
    {是偶函数，它在 \((0,+\infty)\) 上是减函数}
    {是奇函数，它在 \((0,+\infty)\) 上是增函数}
    {是偶函数，它在 \((0,+\infty)\) 上是增函数}

\begin{answer}
C
\end{answer}

\begin{solution}
令 \(f(x)=\frac{\e^x-\e^{-x}}2\). 有 \(f(-x)=-f(x)\)，所以 \(f\) 是奇函数；又
\(f'(x)=\frac{\e^x+\e^{-x}}2>0\)，所以 \(f\) 在 \(\R\) 上严格递增. 同时 \(f(x)\) 在 \(x\to+\infty\) 时趋于 \(+\infty\)，在 \(x\to-\infty\) 时趋于 \(-\infty\)，故它的值域为 \(\R\)，反函数存在.

设反函数为 \(g=f^{-1}\). 由 \(f(-x)=-f(x)\)，对任意 \(y\) 有 \(g(-y)=-g(y)\)，所以 \(g\) 是奇函数；严格递增函数的反函数仍严格递增，因此 \(g\) 在 \((0,+\infty)\) 上是增函数，故选 C.
\end{solution}
\end{problem}

\begin{problem}
如果函数 \(f(x)=x^2+bx+c\) 对任意实数 \(t\) 都有 \(f(2+t)=f(2-t)\)，那么（\quad）

\choices
    {\(f(2)<f(1)<f(4)\)}
    {\(f(1)<f(2)<f(4)\)}
    {\(f(2)<f(4)<f(1)\)}
    {\(f(4)<f(2)<f(1)\)}

\begin{answer}
A
\end{answer}

\begin{solution}
展开两边，得
\(f(2+t)=t^2+(4+b)t+(4+2b+c)\)，
\(f(2-t)=t^2-(4+b)t+(4+2b+c)\). 题设要求它们对任意实数 \(t\) 相等，所以一次项系数必须相等，\(4+b=-(4+b)\)，从而 \(b=-4\). 因此
\(f(x)=x^2-4x+c=(x-2)^2+c\).

于是 \(f(2)=c\)、\(f(1)=c+1\)、\(f(4)=c+4\)，所以 \(f(2)<f(1)<f(4)\)，故选 A.
\end{solution}
\end{problem}

\begin{problem}
长方体的全面积为 \(11\)，\(12\) 条棱长度之和为 \(24\)，则这个长方体的一条对角线长为（\quad）

\choices
    {\(2\sqrt{3}\)}
    {\(\sqrt{14}\)}
    {\(5\)}
    {\(6\)}

\begin{answer}
C
\end{answer}

\begin{solution}
设长方体的三条棱长为 \(a\)，\(b\)，\(c>0\). 全面积为 \(2(ab+bc+ca)=11\)，所以 \(ab+bc+ca=\frac{11}{2}\). 由于每种棱长各有四条，\(12\) 条棱长度之和给出 \(4(a+b+c)=24\)，即 \(a+b+c=6\).

长方体的空间对角线 \(l\) 满足
\(l^2=a^2+b^2+c^2=(a+b+c)^2-2(ab+bc+ca)=6^2-2\cdot\frac{11}{2}=25\). 因为 \(l>0\)，所以 \(l=5\)，故选 C.
\end{solution}
\end{problem}

\section{填空题}

\begin{problem}
方程 \(\frac{1+3^{-x}}{1+3^x}=3\) 的解是\fillinblank{}.

\begin{answer}
\(-1\)
\end{answer}

\begin{solution}
令 \(t=3^x>0\)，则 \(3^{-x}=\frac{1}{t}\)，原方程化为 \(\frac{1+1/t}{1+t}=\frac{1}{t}=3\). 所以 \(t=\frac{1}{3}=3^{-1}\)，从而 \(x=-1\).
\end{solution}
\end{problem}

\begin{problem}
\(\sin 15^{\circ}\sin 75^{\circ}\) 的值是\fillinblank{}.

\begin{answer}
\(\frac{1}{4}\)
\end{answer}

\begin{solution}
由积化和差公式，\(\sin15^{\circ}\sin75^{\circ}=\frac{1}{2}\bigl(\cos60^{\circ}-\cos90^{\circ}\bigr)=\frac{1}{2}\left(\frac{1}{2}-0\right)=\frac{1}{4}\).
\end{solution}
\end{problem}

\begin{problem}
设含有 \(10\) 个元素的集合的全部子集数为 \(S\)，其中由 \(3\) 个元素组成的子集数为 \(T\)，则 \(\frac{T}{S}\) 的值为\fillinblank{}.

\begin{answer}
\(\frac{15}{128}\)
\end{answer}

\begin{solution}
含有 \(10\) 个元素的集合的每个元素都有“选入”或“不选入”两种情况，所以全部子集数为 \(S=2^{10}\). 其中恰含 \(3\) 个元素的子集数为 \(T=\mathrm C_{10}^{3}=\frac{10\cdot9\cdot8}{3\cdot2\cdot1}=120\). 因而 \(\frac{T}{S}=\frac{120}{2^{10}}=\frac{15}{128}\).
\end{solution}
\end{problem}

\begin{problem}
焦点为 \(F_1(-2,0)\) 和 \(F_2(6,0)\)，离心率为 \(2\) 的双曲线的方程是\fillinblank{}.

\begin{answer}
\(\frac{(x-2)^2}{4}-\frac{y^2}{12}=1\)
\end{answer}

\begin{solution}
两焦点的中点为 \((2,0)\)，焦半距为 \(c=4\). 由离心率 \(e=\frac{c}{a}=2\)，得 \(a=2\). 双曲线满足 \(c^2=a^2+b^2\)，所以 \(b^2=16-4=12\). 由于焦点在 \(x\) 轴上，方程为 \(\frac{(x-2)^2}{4}-\frac{y^2}{12}=1\).
\end{solution}
\end{problem}

\begin{problem}
已知等差数列 \(\{a_n\}\) 的公差 \(d\neq 0\)，且 \(a_1\)，\(a_3\)，\(a_9\) 成等比数列，则 \(\frac{a_1+a_3+a_9}{a_2+a_4+a_{10}}\) 的值是\fillinblank{}.

\begin{answer}
\(\frac{13}{16}\)
\end{answer}

\begin{solution}
设 \(a_1=a\)，则 \(a_3=a+2d\)，\(a_9=a+8d\). 三项成等比数列，所以 \((a+2d)^2=a(a+8d)\)，化简得 \(4d(d-a)=0\). 因为 \(d\neq0\)，故 \(a=d\)，从而 \(a_n=nd\). 因此
\(\frac{a_1+a_3+a_9}{a_2+a_4+a_{10}}=\frac{(1+3+9)d}{(2+4+10)d}=\frac{13}{16}\).
\end{solution}
\end{problem}

\section{解答题}

\begin{problem}
已知 \(z\in \mathbf{C}\)，解方程：\(z\bar{z}-3\i\bar{z}=1+3\i\).

\begin{solution}
设 \(z=x+y\i\)，则 \(\bar z=x-y\i\)，且 \(z\bar z=x^2+y^2\). 代入方程并整理实部、虚部，得
\(x^2+y^2-3y-3x\i=1+3\i\). 比较虚部得 \(-3x=3\)，所以 \(x=-1\)；比较实部得 \(1+y^2-3y=1\)，即 \(y(y-3)=0\). 因而 \(y=0\) 或 \(y=3\)，解为 \(z=-1\) 或 \(z=-1+3\i\).
\end{solution}
\end{problem}

\begin{problem}
已知 \(\frac{\pi}{2}<\beta<\alpha<\frac{3\pi}{4}\)，\(\cos(\alpha-\beta)=\frac{12}{13}\)，\(\sin(\alpha+\beta)=-\frac{3}{5}\). 求 \(\sin 2\alpha\) 的值.

\begin{solution}
令 \(\delta=\alpha-\beta\)，则 \(0<\delta<\frac{\pi}{4}\)，所以 \(\sin\delta=\frac{5}{13}\). 又 \(\pi<\alpha+\beta<\frac{3\pi}{2}\)，且 \(\sin(\alpha+\beta)=-\frac{3}{5}\)，故 \(\cos(\alpha+\beta)=-\frac{4}{5}\). 于是
\(\sin2\alpha=\sin\bigl((\alpha+\beta)+(\alpha-\beta)\bigr)=\sin(\alpha+\beta)\cos\delta+\cos(\alpha+\beta)\sin\delta= -\frac{3}{5}\cdot\frac{12}{13}-\frac{4}{5}\cdot\frac{5}{13}=-\frac{56}{65}\).
\end{solution}
\end{problem}

\begin{problem}
已知：两条异面直线 \(a\)，\(b\) 所成的角为 \(\theta\)，它们的公垂线段 \(AA_1\) 的长度为 \(d\). 在直线 \(a\)，\(b\) 上分别取点 \(E\)，\(F\)，设 \(A_1E=m\)，\(AF=n\). 求证：\(EF=\sqrt{d^2+m^2+n^2\pm 2mn\cos\theta}\).

\begin{solution}
由线段 \(A_1E\) 和 \(AF\) 分别沿两直线，取向量分解
\(\overrightarrow{EF}=\overrightarrow{A_1A}+\overrightarrow{AF}-\overrightarrow{A_1E}\). 因为 \(AA_1\) 是两直线的公垂线段，\(\overrightarrow{A_1A}\) 同时垂直于 \(\overrightarrow{A_1E}\) 和 \(\overrightarrow{AF}\)，且 \(|\overrightarrow{A_1A}|=d\)，\(|\overrightarrow{A_1E}|=m\)，\(|\overrightarrow{AF}|=n\). 所以
\(EF^2=d^2+m^2+n^2-2\overrightarrow{AF}\cdot\overrightarrow{A_1E}\). 两直线所成角为 \(\theta\)，而两条有向射线的夹角可能为 \(\theta\) 或 \(\pi-\theta\)，故 \(\overrightarrow{AF}\cdot\overrightarrow{A_1E}=\pm mn\cos\theta\). 于是 \(EF^2=d^2+m^2+n^2\pm2mn\cos\theta\)，从而 \(EF=\sqrt{d^2+m^2+n^2\pm2mn\cos\theta}\).
\end{solution}
\end{problem}

\begin{problem}
设等差数列 \(\{a_n\}\) 的前 \(n\) 项和为 \(S_n\). 已知 \(a_3=12\)，\(S_{12}>0\)，\(S_{13}<0\).

\begin{enumerate}
\item 求公差 \(d\) 的取值范围；
\item 指出 \(S_1\)，\(S_2\)，\(\cdots\)，\(S_{12}\) 中哪一个值最大，并说明理由.
\end{enumerate}

\begin{solution}
由 \(a_3=12\)，得 \(a_1=12-2d\). 因而
\(S_{12}=\frac{12}{2}\bigl(2a_1+11d\bigr)=6(24+7d)>0\)，所以 \(d>-\frac{24}{7}\)；又
\(S_{13}=\frac{13}{2}\bigl(2a_1+12d\bigr)=13(12+4d)<0\)，所以 \(d<-3\). 综合得 \(d\in\left(-\frac{24}{7},-3\right)\).

\(a_n=12+(n-3)d\). 因为 \(d>-\frac{24}{7}\)，有 \(a_6=12+3d>\frac{12}{7}>0\)；因为 \(d<-3\)，有 \(a_7=12+4d<0\). 所以 \(S_n-S_{n-1}=a_n\) 在 \(n=2\)，\(\ldots\)，\(6\) 时为正，在 \(n=7\)，\(\ldots\)，\(12\) 时为负，故 \(S_1\)，\(S_2\)，\(\ldots\)，\(S_{12}\) 先增后减，最大值为 \(S_6\).
\end{solution}
\end{problem}

\begin{problem}
已知椭圆 \(\frac{x^2}{a^2}+\frac{y^2}{b^2}=1\)（\(a>b>0\)），\(A\)，\(B\) 是椭圆上的两点，线段 \(AB\) 的垂直平分线与 \(x\) 轴相交于点 \(P(x_0,0)\). 证明：\(-\frac{a^2-b^2}{a}<x_0<\frac{a^2-b^2}{a}\).

\begin{solution}
设 \(A=(x_1,y_1)\)，\(B=(x_2,y_2)\). 若 \(x_1=x_2\)，则由椭圆方程可知两个不同点只能关于 \(x\) 轴对称，此时 \(AB\) 的垂直平分线就是 \(x\) 轴，不会与其交于唯一的点 \(P\)，故 \(x_1\neq x_2\). 由两点在椭圆上，得
\(y_i^2=b^2\left(1-\frac{x_i^2}{a^2}\right)\)，从而
\(x_i^2+y_i^2=b^2+\frac{a^2-b^2}{a^2}x_i^2\)，\((i=1,2)\).

点 \(P\) 在线段 \(AB\) 的垂直平分线上，所以 \(PA=PB\). 平方并消去 \(x_0^2\)，得
\(2x_0(x_2-x_1)=x_2^2+y_2^2-x_1^2-y_1^2=\frac{a^2-b^2}{a^2}(x_2-x_1)(x_2+x_1)\). 因为 \(x_1\neq x_2\)，所以
\(x_0=\frac{a^2-b^2}{2a^2}(x_1+x_2)\).

又由椭圆方程 \(|x_1|\leqslant a\)，\(|x_2|\leqslant a\)，所以 \(|x_1+x_2|\leqslant2a\)，进而 \(|x_0|\leqslant\frac{a^2-b^2}{a}\). 若取等号，则必须有 \(x_1=x_2=a\) 或 \(x_1=x_2=-a\)，此时 \(A=B\)，与两点为不同点矛盾. 因此严格成立，即
\( -\frac{a^2-b^2}{a}<x_0<\frac{a^2-b^2}{a}\).
\end{solution}
\end{problem}
